\section{Interpretation and Limits}
\label{sec:discussion}

The bank constructs each candidate from the network equations rather than fitting a source-label decision boundary. That is an architectural distinction, not evidence of causal source-excluded generalization. It retains the explanation as states evolve and distinguishes an operator being restored from the grid returning to equilibrium. In this development test the practical gain is timely load attribution; line attribution does not separate the physical and learned references.

That conclusion has three boundaries. First, the causal result uses a known nine-load/nine-line catalog and an intervention-free startup. It is not a test of all 39 buses or arbitrary unknown event families. Second, the instrument and terminal map are frozen research assumptions, including the provisional BUS5 assignment and excluded BUS39 current angle. Agreement within that model does not establish the producer's actual CT placement or exact replication of RAW0001. Third, the network bank is finite and its startup nuisance estimates are plug-in. Its masses are useful attribution weights but are not demonstrated calibrated probabilities.

The nominal reference and failed onset-witness ablation prevent a stronger accuracy claim. More network hypotheses reduce extras in H04 but delay a BUS31 decision; more restrictive admission removes extras but rejects a real BUS16 event. Those are differences in what the decision retains and when it is emitted. Reporting eventual source recovery, or zero extras alone, would select different winners and hide the trade-off.

The mathematical interpretation is equally scoped. Good hidden-voltage reconstruction does not prove event separability. The batch test supports the importance of profiled response overlap, but does not validate every analytic small-amplitude or changing-sensor limit in the supplement. The 137-support batch reference retains its original quadrature; the causal implementation has not been proved to be the unrestricted joint estimator or an exact reduction of it.

The ExtraTrees result tests the particular frozen causal adapters under these physical conditions. Their load failure calls for that restricted interpretation, rather than a general statement that ML cannot recognize loads. Likewise, Chen's colored-noise/DC mismatch is kept visible instead of being used to claim universal superiority over the published algorithm under its own assumptions.

The final decision rests on the prespecified C01 comparison and claim/evidence review, without another bus-specific correction. Faults, generation mechanisms, PMU loss, unseen-source causal transfer and field operation are future scopes, not hidden requirements for reporting this restricted result. The proposed joint estimator remains documented in the supplement so that exploratory work is preserved without being mistaken for the mechanism validated here.
